% ============================================================
% Section 2: Mathematical Foundation
% ============================================================
\section{Mathematical Foundation}

\begin{frame}{The Mathematics (Brief)}
    \textbf{1. Kernel Density Estimation:}

    Let $\mathbf{x} \in \mathbb{R}^d$ and $\mathbf{x}_i \in \mathbb{R}^d$ for $i = 1, 2, \ldots, n$ where $\mathbf{x}_i$ are points within bandwidth $h$ of $\mathbf{x}$.

    \[
    \hat{f}(\mathbf{x}) = \frac{1}{n h^d} \sum_{i=1}^{n} K\left(\frac{\mathbf{x} - \mathbf{x}_i}{h}\right)
    \]

    where:
    \begin{itemize}
        \item $K: \mathbb{R}^+ \to \mathbb{R}$ is a kernel function
        \item $h > 0$ is the bandwidth parameter
        \item $d$ is the dimensionality of the feature space
        \item $\mathbf{x}_i \in \mathbb{R}^d$ are the data points
    \end{itemize}

    \small{\textcolor{customred}{Key insight:} As $h$ increase, we now need to consider more points to compare in the kernel density estimation, which 
    may cause some performance issues.} 
\end{frame}

\begin{frame}{The Mathematics (Contd.)}

    \textbf{2. Mean Shift Vector:}
    \[
    \mathbf{m}(\mathbf{x}) = \left[\frac{\sum_{i=1}^{n} \mathbf{x}_i g\left(\left\|\frac{\mathbf{x} - \mathbf{x}_i}{h}\right\|^2\right)}{\sum_{i=1}^{n} g\left(\left\|\frac{\mathbf{x} - \mathbf{x}_i}{h}\right\|^2\right)}\right] - \mathbf{x}
    \]

    \vspace{0.3cm}

    \textbf{3. Update Rule:}
    \[
    \mathbf{x}^{(t+1)} = \mathbf{x}^{(t)} + \mathbf{m}(\mathbf{x}^{(t)})
    \]

    \small{\textcolor{customred}{Key insight:} Mean shift points toward increasing density}
\end{frame}

\begin{frame}{Common Kernel Functions}
    \begin{columns}[T]
        \begin{column}{0.5\textwidth}
            \textbf{Gaussian Kernel:}
            \[
            K(\mathbf{u}) = \exp\left(-\frac{\|\mathbf{u}\|^2}{2}\right)
            \]

            \vspace{0.5cm}

            \textbf{Epanechnikov Kernel:}
            \[
            K(\mathbf{u}) = \begin{cases}
            1 - \|\mathbf{u}\|^2 & \text{if } \|\mathbf{u}\| \leq 1 \\
            0 & \text{otherwise}
            \end{cases}
            \]
        \end{column}

        \begin{column}{0.5\textwidth}
            \textbf{Flat Kernel:}
            \[
            K(\mathbf{u}) = \begin{cases}
            1 & \text{if } \|\mathbf{u}\| \leq 1 \\
            0 & \text{otherwise}
            \end{cases}
            \]

            \vspace{0.5cm}

            \small{\textcolor{customred}{Choice matters!}}
            \begin{itemize}
                \item Gaussian: smooth, no hard boundary
                \item Epanechnikov: optimal for Mean Shift (it is the most efficient kernel since the derivative is 1 or just 0 (sklearn use this too))
                \item Flat (Uniform kernel): simple, fast computation (but no one use for Mean Shift because the derivative is 0)
            \end{itemize}
        \end{column}
    \end{columns}
\end{frame}


\begin{frame}{Algorithm Steps}
    \begin{algorithm}[H]
        \caption{Mean Shift Clustering}
        \begin{algorithmic}[1]
            \State \textbf{Input:} Data points $\{\mathbf{x}_1, \ldots, \mathbf{x}_n\}$, bandwidth $h$, Kernel function $K$
            \State \textbf{Output:} Cluster assignments
            \State
            \For{each data point $\mathbf{x}_i$}
                \State $\mathbf{y} \gets \mathbf{x}_i$ \Comment{Initialize at data point}
                \While{not converged}
                    \State Compute mean shift: $\mathbf{m}(\mathbf{y})$
                    \State Update: $\mathbf{y} \gets \mathbf{y} + \mathbf{m}(\mathbf{y})$
                \EndWhile
                \State Store converged point (mode) for $\mathbf{x}_i$
            \EndFor
            \State
            \State Group points by their converged modes
            \State \textbf{Return} cluster assignments
        \end{algorithmic}
    \end{algorithm}
\end{frame}
